\begin{question}
\noindent Below is a circle with centre $O$.\\
\noindent $T$ is the point where the tangent line touches the circle.

\begin{center}
\begin{tikzpicture}[scale=2.2]
    % Centre of circle
    \coordinate (O) at (0,0);
    \fill[black] (O) circle (1pt);
    \node[above] at (O) {$O$};

    % Circle (the dish)
    \draw[thick] (O) circle (1);

    % Point of tangency T
    \coordinate (T) at (-90:1);
    \node[below left] at (T) {$T$};

    % Draw radius OT
    \draw[thick] (O) -- (T);

    % Tangent line at T (perpendicular to OT), extended to P
    \coordinate (P) at ($(T) + (2.4,0)$);
    \draw[thick] (T) -- (P) node[right] {$P$};

    % Line from O to P (support cable)
    \draw[thick] (O) -- (P);

    % Mark angle OPT = 24 degrees
    \pic [draw, angle eccentricity=1.4, angle radius=0.75cm, "$24^\circ$"] {angle = O--P--T};

    % Mark angle to find at O (angle TOP = x)
    \pic [draw, angle eccentricity=1.4, angle radius = 0.75cm, "$x^\circ$"] {angle = T--O--P};

\end{tikzpicture}
\end{center}


Angle $OPT = 24^\circ$.

\noindent Work out the size of angle $TOP$, marked $x^\circ$ on the diagram. Give a reason for each stage of your working.

\vspace{3cm}

\begin{flushright}
\answerequnit{x}{$^\circ$}{2}
\end{flushright}
\end{question}
